
\documentclass{beamer}
\setbeamertemplate{navigation symbols}{}
\usepackage[catalan]{babel}
\usepackage[utf8]{inputenc}
\usepackage{ragged2e}
\usepackage{tikz}
\usepackage{amsmath}
\usepackage{amssymb}
\usetheme{Warsaw}
\justifying

\beamertemplatetransparentcovereddynamic
%\beamersetuncovermixins{\opaqueness<1>{25}}{\opaqueness<2->{15}}



\begin{document}
\title{Matemàtiques Primer Batxillerat}
\author{Artur Arroyo}
\date{\today}


\begin{frame}
\titlepage
\end{frame}

\begin{frame}\frametitle{Matemàtiques primer batxillerat}\tableofcontents[pausesubsections]
\end{frame}


\section{Nombres reals}
%\begin{frame}\frametitle{Title}
%Each frame should have a title.
%\end{frame}


\subsection{Classificació dels nombres reals}
\begin{frame}\frametitle{Els nombres naturals}
Comencem recordant els diferents conjunts de nombres dins els
reals.\pause
\begin{block}{Definició}
El conjunt dels nombres \emph{naturals}
$\mathbb{N}=\{1,2,3,\ldots\}$
\end{block}
\pause
\begin{exampleblock}{Exemples}
$$123\in \mathbb{N},\quad -5\notin\mathbb{N}, \quad
\frac{2}{9}\notin\mathbb{N}, \quad \sqrt{2}\notin\mathbb{N}, \quad
\pi\notin\mathbb{N}$$
\end{exampleblock}\pause
\begin{alertblock}{Atenció}
$$0\notin\mathbb{N}$$
\end{alertblock}
\end{frame}
\begin{frame}\frametitle{Els nombres enters}
Si volem incloure el zero hem d'ampliar el conjunt dels nombres
naturals. Això es fa considerant els anomenats nombres
enters\pause
\begin{block}{Definició}
El conjunt dels nombres \emph{enters} $\mathbb{Z}=\{\ldots,
-3,-2,-1,0, 1,2,3,\ldots\}$
\end{block}\pause
\begin{exampleblock}{Exemples}
$$123\in \mathbb{Z},\quad -5\in\mathbb{Z}, \quad
\frac{2}{9}\notin\mathbb{Z}, \quad \sqrt{2}\notin\mathbb{Z}, \quad
\pi\notin\mathbb{Z}$$
\end{exampleblock}
\end{frame}
\begin{frame}
\begin{alertblock}{Atenció!}
Sovint es parla de nombres positius, negatius, i hi pot haver una
mica de confusió. Hem de tenir clar que els naturals $\mathbb{N}$
són els enters $\mathbb{Z}$ que són \emph{positius} i d'aquesta
forma veiem que quan es parla de nombres positius el zero no hi
és. \\
Si volem considerar $\mathbb{N}\cup \{0\}$, haurem
d'anomenar aquests nombres \emph{no negatius}.
\end{alertblock}
\end{frame}

\begin{frame}\frametitle{Els nombres racionals}
És clar que nombres com $\frac{3}{7}$ ó $\frac{-5}{11}$ no són
enters a pesar de que el numerador i el denominador sí que ho són.
Per poder treballar amb les fraccions cal estendre els
enters.\pause
\begin{block}{Definició}
El conjunt dels nombres \emph{racionals}
$\mathbb{Q}=\{\frac{a}{b},$ amb $a, b \in \mathbb{Z}\}$
\end{block}\pause
\begin{exampleblock}{Exemples}
$$123\in \mathbb{Q},\quad \frac{2}{9}\in\mathbb{Q}, \quad \sqrt{2}\notin\mathbb{Q}, \quad
\pi\notin\mathbb{Q}$$
\end{exampleblock}
\end{frame}

\begin{frame}\frametitle{Fracció generatriu d'un nombre decimal}
 Arribats a aquest punt hem de pensar com tractar
nombres com $1,23\quad 2,7272\ldots\quad 4,9666\ldots$ Aquests
nombres són racionals, (això s'ha de justificar), i sempre es pot
trobar l'anomenada \emph{fracció generatriu}. Perquè podem afirmar
que són racionals? Fixem-nos que el primer té dues xifres
decimals, al segon, malgrat que en té un nombre infinit, sempre es
repeteix el mateix patró i el tercer també té un patró que es
repeteix a pesar de que hi ha xifres diferents entre ell i la
coma.
\end{frame}


\begin{frame}\frametitle{Fracció generatriu d'un nombre decimal}
\begin{block}{Nombres decimals. Classificació}

$\text{Nombres decimals}\begin{cases} \text{Decimals periòdics}&
\begin{cases}
\text{Periòdics purs}\\
\text{Periòdics mixtos}
\end{cases}\\
\text{Decimals exactes}
\end{cases}$
\end{block}\pause
\begin{exampleblock}{Exemples}
El nombre 1,23 és decimal exacte. $2,7272\ldots$ és decimal
periòdic pur (el periode és 72) i $4,9666\ldots$ és decimal
periòdic mixt (el periode és 6)
\end{exampleblock}
\end{frame}
\begin{frame}\frametitle{Fracció generatriu d'un nombre decimal exacte}
\begin{itemize}
\item
El nombre 1,23 és decimal exacte. La seva fracció generatriu es
troba de forma immediata,\\ $$1,23=\frac{123}{100}$$
\end{itemize}
\end{frame}
\begin{frame}\frametitle{Fracció generatriu d'un nombre decimal periòdic pur}
\begin{itemize}
\item
El nombre $2,7272\ldots$ és decimal periòdic pur, per trobar la
seva fracció genetratriu seguim el següent procés:
\end{itemize}\pause
\begin{itemize}
\item[]
\begin{align}
 \text{anomenem x al nombre decimal,}  & &x&=2,7272\ldots \\\pause
\text{multipliquem per 100,}&&100x&= 272,7272\ldots \\\pause
\text{restem (2) de (1),}& &99x&=270 \\\pause \text{aïllant la x,}
& &x &=\frac{270}{99}=\frac{30}{11}
\end{align}
%\begin{eqnarray*}
 %0 & = & x^2 -3x -1  \\ \pause
 %0 & = &  x^2 -3x + (-3/2)^2 - (3/2)^2 -1  \\ \pause
 %0 & = &  (x - 3/2)^2  -9/4 - 4/4  \\ \pause
 %0 & = &   (x - 3/2)^2  -13/4
%\end{eqnarray*}
\end{itemize}\pause
Es comprova fàcilment amb qualsevol calculadora que
$\frac{30}{11}=2,7272\ldots$ de forma que aquest nombre decimal és
racional.
\end{frame}

\begin{frame}\frametitle{Fracció generatriu d'un nombre decimal periòdic mixte}
\begin{itemize}
\item
El nombre $4,9666\ldots$ és decimal periòdic mixte, per trobar la
seva fracció genetratriu seguim el següent procés:\pause
\end{itemize}
\begin{itemize}
\item[]
\begin{align}
 \text{anomenem x al nombre decimal,}  & &x&=4,9666\ldots \\\pause
\text{multipliquem per 100,}&&100x&= 496,666\ldots \\\pause
\text{ara multipliquem (5) per 10,}&&10x&= 49,666\ldots \\\pause
\text{restem (7) de (6),}& &90x&=447 \\\pause \text{aïllant la x,}
& &x &=\frac{447}{90}=\frac{149}{30}
\end{align}
\end{itemize}
\end{frame}

\begin{frame}\frametitle{Propietat de densitat dels racionals}
\begin{alertblock}{Important!}
Entre dos nombres racionals qualssevol $a, b$ sempre se'n pot
trobar un altre, n'hi ha prou de fer l'operació $\frac{a+b}{2}$.
Aquesta propietat és prou important com per a tenir nom i es farà
servir més endevant.
\end{alertblock}
\end{frame}
\begin{frame}
\begin{block}{Definició}
El conjunt dels nombres racionals $\mathbb{Q}$ és \emph{dens} a
$\mathbb{R}$.
\end{block}\pause
\begin{exampleblock}{Exemple 1}
Quin nombre és just al mig de 3 i 11? Resposta:
$\frac{3+11}{2}=\frac{14}{2}=7$
\end{exampleblock}\pause
\begin{exampleblock}{Exemple 2}
Quin nombre és just al mig de $\frac{3}{5}$ i $\frac{11}{7}$?\\
Resposta:
$\frac{\frac{3}{5}+\frac{11}{7}}{2}=\frac{\frac{21+55}{35}}{2}=\frac{\frac{76}{35}}{2}=\frac{76}{70}=\frac{28}{35}$

\end{exampleblock}
\end{frame}
\begin{frame}\frametitle{Operacions amb fraccions}
Als exemples anteriors hi ha un parell d'operacions que, tot i ser
molt bàsiques, s'han de tenir molt clares ja que s'utilitzaran
sovint al llarg del curs.\pause
\begin{block}{Suma de nombres racionals}
Per a sumar dues fraccions $\frac{A}{B}$, $\frac{C}{D}$ sempre
podem fer $$\frac{A}{B}+\frac{C}{D}=\frac{AD+BC}{BD}$$
\end{block}\pause
\small A vegades pot ser més ràpid reduir a comú denominador les
fraccions. En cada cas s'haurà de valorar quin mètode fer
servir.\pause

\normalsize
\begin{block}{Fraccions amb fraccions}
Farem servir la fòrmula:
$\frac{\frac{A}{B}}{\frac{C}{D}}=\frac{AD}{BC}$
\end{block}
\end{frame}
\begin{frame}\frametitle{Els nombres irracionals $\mathbb{I}$}
Hi ha nombres que tenen infinites xifres decimals però al no ser
periòdics no se'n pot trobar la fracció generatriu i per tant, no
són racionals.\pause
\begin{exampleblock}{Exemples}
$$\sqrt{3}\in\mathbb{I},\quad \sqrt[3]{7}\in\mathbb{I},\quad \pi\in\mathbb{I},\quad
e\in\mathbb{I},\quad \phi\in\mathbb{I}$$
\end{exampleblock}\pause
Noteu que no tots els nombres irracionals són arrels d'algun
nombre. Entre els nombres irracionals hi ha les constants notables
més importants. A l'exemple apareixen el nombre pi ($\pi$), el
nombre d'Euler ($e$) i l'anomenada raó àurea ($\phi$).\pause
\begin{alertblock}{Compte!}
$$\sqrt{4}=2, \quad \sqrt[3]{-125}=-5,\quad \sqrt{-1}\notin\mathbb{R},\ldots$$ No totes les arrels són
nombres irracionals!
\end{alertblock}
\end{frame}
\begin{frame}{Els nombres reals}
Amb tot el que hem vist abans podem establir la següent relació
d'inclusió entre els diferents conjunts.
\begin{block}{Els nombres reals}
$$\mathbb{N}\subset\mathbb{Z}\subset\mathbb{Q}\subset\mathbb{R}$$
$$\mathbb{Q}\cup\mathbb{I}=\mathbb{R}$$
$$\mathbb{Q}\cap\mathbb{I}=\{\emptyset\}$$
\end{block}
\end{frame}

\subsection{La recta real}
\begin{frame}\frametitle{Relació d'ordre}
\begin{itemize}
\item
 Els nombres reals es representen còmodament sobre una recta
amb uns símbols a dreta ($\infty$) i esquerra (-$\infty$) que
completen els reals, en el sentit que no hi ha cap nombre real més
gran que infinit.\pause
\begin{alertblock}{Compte!}
Aquests símbols \emph{no} són nombres reals:
$\quad\pm\infty\notin\mathbb{R}$
\end{alertblock}\pause
\item
 En el conjunt dels nombres reals s'estableix
una relació d'ordre que ens permet decidir si un nombre és més
gran, més petit o igual que un altre.\pause
\begin{block}{Relació d'ordre en els nombres reals}
Direm que un nombre $a$ és més gran que un altre $b$, si $a$ està
a la dreta que $b$ a la recta real.
\end{block}
\end{itemize}
\end{frame}

\begin{frame}\frametitle{Intervals}
A la recta real es poden considerar diferents tipus de conjunts,
anome\-nats \emph{intervals}. Segons incloguin o no els seus
extrems es classifiquen en:\pause
%\begin{block}{Intervals oberts}
%Si un conjunt o interval no conté els seus extrems l'anomerarem
%\emph{obert}.
%\end{block}
%
%\begin{exampleblock}{Exemple d'interval obert}
%Sigui el conjunt $A=(-3, 10)$. Llavors: $-4\notin A$, $-3\notin
%A$, $4\in A$,
%\end{exampleblock}
\begin{itemize}
\item
Interval obert: $(a,b)$. Són els nombres estrictament més grans
que $a$ i estrictament més petits que $b$.\pause
\item
Interval tancat: $[a,b]$. Són els nombres més grans o iguals que
$a$ i més petits o iguals que $b$.\pause

\item
Interval semiobert: $[a,b), \; (a,b]$. El primer conté $a$ però no
$b$ i el segon a l'inrevés .
\end{itemize}\pause
\emph{També es poden considerar intervals infinits}.
\end{frame}
\begin{frame}
\begin{exampleblock}{Exemples}
Si considerem els conjunts: $$A=(-3, 10), \;B=[-5,-1), \;
C=(2,\infty), \; D=(-\infty,12]$$\\
$\;$\\
 Llavors: $$\quad-4\notin A,\quad
-3\notin A,\quad 9\in A,\quad 10\notin A,\quad 11\notin A,\quad
-5\in B$$\\
$$2\notin C,\quad 0\in C,\quad 12\in D,\quad 12,0001\notin D$$
\end{exampleblock}\pause
Noteu que, en qualsevol cas, $\infty$ o $-\infty$ sempre es prenen
oberts, ja que no són nombres reals.
\end{frame}

\begin{frame}\frametitle{Notació científica}
\begin{itemize}
\item Un nombre en notació científica és de la forma
$\mathbf{a\cdot10^{b}}$ on $|a|$ és un nombre decimal exacte de
l'interval $[1,10)$ i l'exponent, $b$, és un nombre enter.
Anomenarem el nombre $a$ \emph{mantissa} i el nombre $b$
\emph{ordre de magnitud}.\pause
\item
La notació científica es fa servir per representar nombres molt
grans o molt petits.
\end{itemize}\pause
\begin{exampleblock}{Exemples}
Moltes constants físiques notables s'escriuen en notació
científica:
\begin{itemize}
\item
velocitat de la llum en el buit, $c=3\cdot10^{8}m/s$
\item
nombre d'Avogadro, $N_{A}=6,02\cdot10^{23} mol^{-1}$
\item
càrrega de l'electró, $q_{e}=1,62\cdot10^{-19}C$
\end{itemize}
\end{exampleblock}
\end{frame}
\subsection{Radicals, operacions i racionalització}
\begin{frame}\frametitle{Radicals}
Recordem amb alguns exemples, com es calculen les potències.\pause
\begin{itemize}
\item
Com es calcula $3^{5}$ ?\pause
$$3^{5}=3\cdot3\cdot3\cdot3\cdot3$$\pause
\item
Com es calcula $3^{-4}$ ?\pause
$$3^{-4}=\frac{1}{3^{4}}$$\pause
\item
Com es calcula $3^{\frac{1}{2}}$ ? i $3^{\frac{5}{7}}$ ?\\
\end{itemize}
\end{frame}
\begin{frame}\frametitle{Radicals}
Per representar aquests nombres farem servir el símbol de l'arrel,
de forma que les arrels no són més que nombres amb exponent
fraccionari.\pause
\begin{block}{Arrel d'un nombre}
Definim $\qquad a^{\frac{m}{n}}=\sqrt[n]{a^{m}}$
\end{block}\pause
Llavors tenim:
\begin{exampleblock}{Exemples}
$3^{\frac{1}{2}}=\sqrt{3}$ i $3^{\frac{5}{7}}=\sqrt[7]{5}$
\end{exampleblock}
\end{frame}
\begin{frame}\frametitle{Operacions amb radicals}
\begin{itemize}
\item
\textbf{Reduir radicals a índex comú.}\\
Farem servir les propietats que ja coneixem per reduir fraccions a
comú denominador.\pause
\begin{exampleblock}{Exemple}
$$\sqrt{2}\; ; \sqrt[3]{3}\; ;\sqrt[5]{5}$$\\
Expressem els radicals com a potències d'exponent fraccionari:
$$2^{\frac{1}{2}},\; 3^{\frac{1}{3}},\; 5^{\frac{1}{5}}$$
Tenim que m.c.m.(2,3,5)=30. Reduïm a comú denominador els
exponents i tornem a expressar els nombres com a radicals:
$$2^{\frac{1}{2}}=2^{\frac{15}{30}}$$
\end{exampleblock}
\end{itemize}
\end{frame}
\begin{frame}
\begin{itemize}
\item
\textbf{Treure o introduir factors de dins de l'arrel}\\
Pot ser útil per simplificar expressions.\pause
\begin{exampleblock}{Exemple}
$\sqrt[3]{a^{4}b^{2}c^{9}}=ac^{3}\sqrt[3]{ab^{2}}\qquad
ab^{2}c^{3}\sqrt[5]{a^{3}bc}=\sqrt[5]{a^{8}b^{11}c^{16}}$
\end{exampleblock}\pause
\item
\textbf{Sumar o restar radicals}
\begin{exampleblock}{Exemple}
\begin{align*}
\sqrt{20}-3\sqrt{125}+2\sqrt{45}=&\\
&=\sqrt{2^{2}5}-3\sqrt{5^{3}}+2\sqrt{3^{2}5}\\
&=2\sqrt{5}-3\cdot5\sqrt{5}+2\cdot3\sqrt{5}\\
&=-7\sqrt{5}
\end{align*}
\end{exampleblock}
\end{itemize}
\end{frame}
\begin{frame}
\begin{itemize}
\item
\textbf{Multiplicar o dividir radicals}\\
Sovint caldrà reduir a índex comú.\pause
\begin{exampleblock}{Exemple}
$\frac{\sqrt[3]{5}\cdot\sqrt{4}}{\sqrt[4]{3}}=\frac{\sqrt[6]{5^{2}}\cdot\sqrt[6]{4^{3}}}{\sqrt[4]{3}}=$
\end{exampleblock}\pause
\item
\textbf{Potències de radicals}\\
En aquest cas n'hi ha prou d'expressar els radicals en forma de
potència i aplicar propietats ja conegudes.\pause

\begin{exampleblock}{Exemple}
$\sqrt[3]{\sqrt{\frac{a^{12}}{a^{18}}}}=
\left(\left(\frac{a^{12}}{a^{18}}\right)^{\frac{1}{2}}\right)^{\frac{1}{3}}$
\end{exampleblock}
\end{itemize}
\end{frame}
\begin{frame}\frametitle{Racionalització}
Racionalitzar consisteix a ''treure'' les arrels que hi pugui
haver al denominador.\pause
\begin{block}{Fraccions del tipus $\frac{a}{\sqrt[n]{b}}$}
En aquest cas n'hi ha prou de multiplicar i dividir per
$\sqrt[n]{b^{n-1}}$
\end{block}\pause
\begin{exampleblock}{Exemple}
$\frac{6}{\sqrt[7]{3^{2}}}=\frac{6}{\sqrt[7]{3^{2}}}\cdot\frac{\sqrt[7]{3^{5}}}{\sqrt[7]{3^{5}}}=\frac{6\cdot\sqrt[7]{3^{5}}}{3}$
\end{exampleblock}
\end{frame}

\begin{frame}\frametitle{Racionalització}
\begin{block}{Fraccions del tipus $\frac{a}{b\pm\sqrt{c}},\; \frac{a}{\sqrt{b}-\sqrt{c}}$}
En aquest cas n'hi ha prou de multiplicar i dividir pel conjugat
del denominador
\end{block}\pause
\begin{exampleblock}{Exemple}
$\frac{4}{\sqrt{3}+\sqrt{5}}=\frac{4}{\sqrt{3}+\sqrt{5}}\cdot\frac{\sqrt{3}-\sqrt{5}}{\sqrt{3}-\sqrt{5}}$
\end{exampleblock}
\end{frame}

\subsection{Logaritmes}
\begin{frame}\frametitle{Logaritmes}
Per tal d'entendre què és el logaritme comencem amb unes
exemples.\pause
\begin{exampleblock}{Exemples}
\begin{itemize}
\item
A quin nombre hem d'elevar 2 per tal que doni 16?\\
\pause $2^{?}=16$\\ \pause
$2^{4}=16$\\
\pause La resposta és 4.\pause
\item
A quin nombre hem d'elevar 10 per tal que doni 1000?\\
\pause $10^{?}=1000$\\
\pause La resposta és 3.\pause
\item
A quin nombre hem d'elevar 7 per tal que doni 49?\\
\pause La
resposta és 2.
\end{itemize}
\end{exampleblock}
\end{frame}
\begin{frame}\frametitle{Logaritmes}
Les preguntes anteriors es poden reecriure mitjançant
logaritmes:\pause
\begin{itemize}
\item
A quin nombre hem d'elevar 2 per tal que doni 16? (Resposta:
4)\\\pause És equivalent a preguntar, quant val el logaritme de 16
en base 2? ($\log_{2}16=4$)\pause
\item
A quin nombre hem d'elevar 10 per tal que doni 1000?\\\pause Quant
val el logaritme de 1000 en base 10? ($\log_{10}1000=3$)\pause
\item
A quin nombre hem d'elevar 7 per tal que doni 49?\\\pause Quant
val el logaritme de 49 en base 7? ($\log_{7}49=2$)
\end{itemize}
\end{frame}
\begin{frame}\frametitle{Definició de logaritme}
Ara ja estem en condicions de donar la definició de logaritme
\begin{block}{Definició}
$$\log_{b}a=c\Leftrightarrow b=a^{c}$$
\end{block}
\pause El nombre $b$, s'anomena \emph{base} del logaritme i el
nombre $a$, s'anomena \emph{argument} del logaritme.

\end{frame}
\begin{frame}
\begin{alertblock}{Atenció!}
Quan la base del logaritme és 10, no cal escriure-la
\end{alertblock}\pause
\begin{exampleblock}{Exemple}
$\log_{10}73\equiv\log73$
\end{exampleblock}\pause
\begin{alertblock}{Atenció}
Quan la base del logaritme és la constant d'Euler, $e$, llavors el
logaritme s'anomena natural i s'escriu de forma lleugerament
diferent
\end{alertblock}\pause
\begin{exampleblock}{Exemple}
$\log_{e}73\equiv\ln73$
\end{exampleblock}
\end{frame}
\begin{frame}\frametitle{Propietats dels logaritmes}
\begin{itemize}
\item
Casos particulars:\\
$$\log_{a}1=0,\qquad \log_{a}a=1$$\pause
%\end{itemize}
\item
Propietats importants
\begin{itemize}
\item
$\log_{b}ac=\log_{b}a+\log_{b}c$\pause
\item
$\log_{b}\left(\frac{a}{c}\right)=\log_{b}a-\log_{b}c$\pause
\item
$\log_{b}a^{n}=n\log_{b}a$
\end{itemize}
\end{itemize}
\end{frame}
\begin{frame}
\begin{exampleblock}{Exemple}
Fent servir les propietats dels logaritmes, calcular:
$$\log_{5}625-\log_{9}81+\log_{64}8$$\pause
Solució:
\begin{align*}
\log_{5}625-\log_{9}81+\log_{64}8=&&\\\pause
\log_{5}5^{4}-\log_{9}9^{2}+\log_{64}64^{\frac{1}{2}}=&\\\pause
4\log_{5}5-2\log_{9}9+\frac{1}{2}\log_{64}64= 4-2+\frac{1}{2}=\frac{5}{2}\\
\end{align*}
\end{exampleblock}
\end{frame}



\end{document}
